\chapter{Caching, Routing und Guardrails}
\label{chap:caching}

% ============================================================
% AUTOR-NOTIZ: Caching/Routing Story (NEU)
% ============================================================
\autornotiz{
2024: SupportPilot FAQ-Bot. 10k Requests/Tag, 60\% wiederholte Fragen (Status, Rückgabe, Versand).
Erster Versuch: Kein Cache. Jeder Request geht an LLM. Kosten \$180/Tag, p99 1.8s.
Nach Caching: Redis + Semantic Cache (text-embedding-3-small, Threshold 0.92).
Hit Rate 38\%, p99 620ms, Kosten \$112/Tag.
Routing: Intent-Classifier (MiniLM) routet Order→gpt-4o-mini, FAQ→cached, Complaint→gpt-4o.
Ergebnis: 62\% Kostenersparnis, p99 unter 1s.
Der Unterschied war nicht der Cache. Es war Intent-Routing + Cache-Keys mit Tenant-Isolation.
}

% ============================================================
% MOTIVATION
% ============================================================
\section{Motivation}

Ein unoptimiertes LLM-System ruft bei jeder Anfrage das Modell auf -- unabhängig davon, ob dieselbe oder eine ähnliche Frage bereits beantwortet wurde. In einem typischen SaaS-System sind 20--40\,\% aller Anfragen Duplikate oder semantische Varianten: \glqq Wie ist mein Bestellstatus?\grqq vs \glqq Wo bleibt meine Bestellung?\grqq.

Jeder vermeidbare API-Call kostet Geld, frisst Latenz und erzeugt CO\textsubscript{2}. Caching eliminiert diese Overheads. Aber LLM-Caching ist komplexer als klassisches HTTP-Caching: Dieselbe Frage kann unterschiedlich beantwortet werden, ähnliche Fragen brauchen ähnliche Antworten, und nicht jede Antwort darf gecached werden.

Parallel zum Caching braucht jedes Produktionssystem \textbf{Guardrails} -- Validierungsschichten, die schädlichen Input filtern, sensitive Daten schützen und die Ausgabe auf Konformität prüfen. Ohne Guardrails ist jedes LLM-System ein Sicherheitsrisiko.

Und zwischen diesen Ebenen vermittelt der \textbf{Semantic Router}: Er entscheidet anhand der Intent-Erkennung, welches Modell, welche Pipeline oder welcher Cache für eine Anfrage zuständig ist.

% ============================================================
% GRUNDLAGEN
% ============================================================
\section{Grundlagen -- Die drei Caching-Ebenen}

LLM-Caching operiert auf drei Ebenen, die sich ergänzen:

\begin{verbatim}
   +------------------------------------------------------------------+
   |                      LLM-Caching-Stack                            |
   |                                                                   |
   |  Ebene 1: Prompt Caching (KV-Cache)                              |
   |  Gleicher Prompt-Praefix -> KV-Cache wiederverwenden             |
   |  Spart: Prefill-Zeit (50-90%), reduziert Latenz                 |
   |  Provider: Anthropic, OpenAI, Google, vLLM                       |
   |                                                                   |
   |  Ebene 2: Response Caching (exakter Match)                       |
   |  Gleicher Input -> gleiche Antwort aus Cache                     |
   |  Spart: Kompletten LLM-Call (100%)                               |
   |  Einsatz: FAQ, Statusabfragen, sich wiederholende Queries        |
   |                                                                   |
   |  Ebene 3: Semantic Caching (Embedding-basiert)                   |
   |  Ahnlicher Input -> ahnliche Antwort aus Cache                   |
   |  Spart: 20-40% der API-Calls in Support-Szenarien                |
   |  Einsatz: Kunden-Support, Dokumentation, interne Wikis           |
   +------------------------------------------------------------------+
\end{verbatim}

\subsection{Prompt Caching (KV-Cache)}

Der effektivste Caching-Mechanismus ist der KV-Cache. Jeder LLM-Call berechnet während der Prefill-Phase Key- und Value-Vektoren für alle Input-Token. Wenn ein Teil des Inputs exakt gleich bleibt (z.\,B. der System-Prompt), kann dieser Teil des KV-Cache wiederverwendet werden:

\begin{lstlisting}[language=Python, caption=Prompt Caching mit Anthropic und OpenAI]
# Anthropic: Prompt Caching via cache_control
response = client.messages.create(
    model="claude-sonnet-4-20250514",
    max_tokens=1024,
    system=[
        {
            "type": "text",
            "text": SYSTEM_PROMPT,   # 2.000 Tokens, statisch
            "cache_control": {"type": "ephemeral"},
        }
    ],
    messages=[{"role": "user", "content": user_query}],
)
print(f"Cache: {response.usage.cache_creation_input_tokens > 0}")

# OpenAI: Prompt Caching (automatisch ab 1.024 Tokens Prefix)
response = client.chat.completions.create(
    model="gpt-4o",
    messages=[
        {"role": "system", "content": SYSTEM_PROMPT},
        {"role": "user", "content": user_query},
    ],
)
print(f"Cache Hit: {response.usage.prompt_tokens_details.cached_tokens > 0}")
\end{lstlisting}

Prompt Caching spart 50--90\,\% der Prefill-Zeit und bis zu 50\,\% der Input-Kosten (Provider berechnen gecachte Tokens günstiger).

\subsection{Response Caching}

Der einfachste Cache: Exakt gleicher Input bekommt die exakt gleiche Antwort:

\begin{lstlisting}[language=Python, caption=Response Cache mit Hash-Key]
import hashlib
import json
from diskcache import Cache  # pip install diskcache

class ResponseCache:
    def __init__(self, cache_dir: str = ".llm_cache",
                 ttl: int = 3600):
        self.cache = Cache(cache_dir)
        self.ttl = ttl

    def _key(self, messages: list[dict], model: str) -> str:
        content = json.dumps(messages, sort_keys=True) + model
        return hashlib.sha256(content.encode()).hexdigest()

    def get(self, messages: list[dict], model: str) -> str | None:
        key = self._key(messages, model)
        return self.cache.get(key)

    def set(self, messages: list[dict], model: str, response: str):
        key = self._key(messages, model)
        self.cache.set(key, response, expire=self.ttl)

    def get_or_call(self, messages, model, llm_func):
        cached = self.get(messages, model)
        if cached:
            return cached, True  # Hit
        response = llm_func(messages, model)
        self.set(messages, model, response)
        return response, False  # Miss

# Nutzung
cache = ResponseCache(ttl=3600)

def answer(messages, model):
    response, is_cached = cache.get_or_call(
        messages, model, lambda m, md: call_llm(m, md)
    )
    status = "(CACHED)" if is_cached else "(GENERATED)"
    print(f"{status} Score: {cache.lookup(messages)[1]:.2f}")
    return response
\end{lstlisting}

\textbf{Einsatzbereiche:} FAQ-Systeme, Statusabfragen, wiederkehrende Extraktionen, Template-basierte Responses.

\praxishinweis{Response Cache ist der einfachste Hebel. SHA-256-Hash über Messages + Model -- und 20-40\% der API-Calls verschwinden. Aber Vorsicht: personalisierte Antworten nie cachen, TTL immer setzen.}

% ===========================================================%
% SEMANTIC CACHING
% ===========================================================%
\section{Semantisches Caching}

Der mächtigste -- und komplexeste -- Cache-Ansatz: Ähnliche Anfragen erkennen und die zwischengespeicherte Antwort einer semantisch äquivalenten Frage zurückgeben.

\subsection{Architektur}

\begin{verbatim}
   [User Query] --> [Embedding-Modell] --> [Vector Search]
                                               |
                                    +----------+----------+
                                    |                     |
                               [Hit > Threshold]    [Miss < Threshold]
                                    |                     |
                               [Cache Response]    [LLM Call]
                                                       |
                                               [Store in Cache]
\end{verbatim}

\subsection{Implementierung}

\begin{lstlisting}[language=Python, caption=Semantic Cache mit Embedding-Vergleich]
import numpy as np
from openai import OpenAI

class SemanticCache:
    def __init__(self, similarity_threshold: float = 0.92):
        self.threshold = similarity_threshold
        self.embeddings: list[np.ndarray] = []
        self.responses: list[str] = []
        self.queries: list[str] = []
        self.client = OpenAI()

    def _embed(self, text: str) -> np.ndarray:
        response = self.client.embeddings.create(
            model="text-embedding-3-small",
            input=text,
        )
        return np.array(response.data[0].embedding)

    def _cosine_similarity(self, a: np.ndarray,
                           b: np.ndarray) -> float:
        return np.dot(a, b) / (np.linalg.norm(a) * np.linalg.norm(b))

    def lookup(self, query: str) -> tuple[str | None, float]:
        query_emb = self._embed(query)

        best_score = 0.0
        best_idx = -1

        for idx, stored_emb in enumerate(self.embeddings):
            score = self._cosine_similarity(query_emb, stored_emb)
            if score > best_score:
                best_score = score
                best_idx = idx

        if best_score >= self.threshold and best_idx >= 0:
            return self.responses[best_idx], best_score
        return None, best_score

    def store(self, query: str, response: str):
        emb = self._embed(query)
        self.embeddings.append(emb)
        self.responses.append(response)
        self.queries.append(query)

    def get_or_generate(self, query: str,
                        llm_func) -> tuple[str, bool]:
        cached, score = self.lookup(query)
        if cached:
            return cached, True
        response = llm_func(query)
        self.store(query, response)
        return response, False

# Nutzung
cache = SemanticCache(threshold=0.92)

def answer(query):
    response, is_cached = cache.get_or_generate(
        query, lambda q: call_llm(q)
    )
    status = "(CACHED)" if is_cached else "(GENERATED)"
    print(f"{status} Score: {cache.lookup(query)[1]:.2f}")
    return response
\end{lstlisting}

\subsubsection{Threshold-Wahl}

Der Similarity-Threshold bestimmt das Verhältnis von Cache-Hits zu Korrektheit:

\begin{table}[H]
\centering
\begin{tabularx}{\linewidth}{lXX}
\toprule
\textbf{Threshold} & \textbf{Hit-Rate} & \textbf{Risiko} \\
\midrule
0,95 & 10--15\,\% & Sehr konservativ, kaum false positives \\
0,92 & 20--30\,\% & Guter Kompromiss für Support-Szenarien \\
0,88 & 30--40\,\% & Aggressiv, false positives möglich \\
0,85 & 40--50\,\% & Nur bei sehr ähnlichen Queries (Templates) \\
\bottomrule
\end{tabularx}
\caption{Einfluss des Thresholds auf Cache-Hits und Korrektheit}
\label{tab:threshold}
\end{table}

\textbf{Praxis:} Starte mit 0,92 für Support-Systeme, 0,88 für FAQ-Systeme. Bei false positives den Threshold erhöhen.

\praxishinweis{Semantischer Cache braucht dynamischen Threshold. Medizinische Fragen: 0.99 (fast Exact Match). Smalltalk: 0.85 reicht. Und immer TTL setzen -- ein gecachtes Modell-Update vom Januar darf nicht im Juni ausgeliefert werden.}

% ===========================================================%
% SEMANTIC ROUTER
% ===========================================================%
\section{Semantic Router}

Ein Semantic Router klassifiziert jede Anfrage nach Intent und routed sie an das optimale Modell oder die optimale Pipeline:

\begin{verbatim}
   [User Query] --> [Intent-Classifier]
                          |
           +--------------+--------------+
           |              |              |
      [Simple Q]    [Complex Q]    [Code Q]
           |              |              |
      [GPT-4o-mini]  [Claude Opus]  [Codestral]
           |              |              |
           +------+-------+------+------+
                  |              |
            [Cache Check]  [Refusal Check]
                  |              |
             [Response]     [Fallback]
\end{verbatim}

\begin{lstlisting}[language=Python, caption=Semantic Router mit Intent-Klassifikation]
from enum import Enum

class Intent(Enum):
    SIMPLE = "simple"
    COMPLEX = "complex"
    CODE = "code"
    REFUSAL = "refusal"  # Unerlaubte Anfragen

class SemanticRouter:
    def __init__(self):
        self.routes = {
            Intent.SIMPLE: {
                "model": "gpt-4o-mini",
                "max_tokens": 256,
                "cost_per_call": 0.0001,
            },
            Intent.COMPLEX: {
                "model": "claude-sonnet-4",
                "max_tokens": 2048,
                "cost_per_call": 0.003,
            },
            Intent.CODE: {
                "model": "codestral-latest",
                "max_tokens": 1024,
                "cost_per_call": 0.002,
            },
        }
        self.classifier_model = "gpt-4o-mini"

    def classify(self, query: str) -> Intent:
        prompt = (
            "Klassifiziere die folgende Anfrage in eine Kategorie.\n"
            f"- '{Intent.SIMPLE.value}': Einfache Frage, "
            f"z.B. Status, Datum, Faktenabfrage\n"
            f"- '{Intent.COMPLEX.value}': Komplexe Analyse, "
            f"mehrschrittiges Reasoning, Vergleich\n"
            f"- '{Intent.CODE.value}': Code-Generierung, "
            f"Debugging, SQL, Regex\n"
            f"- '{Intent.REFUSAL.value}': Illegal, gefaehrlich, "
            f"personenbezogene Daten, Passwoerter\n"
            f"\nAnfrage: {query}\n"
            f"Antwort NUR mit dem Kategorie-String."
        )
        response = call_llm(self.classifier_model, prompt,
                            temperature=0.0)
        try:
            return Intent(response.strip().lower())
        except ValueError:
            return Intent.SIMPLE

    def route(self, query: str, user_id: str = None) -> dict:
        intent = self.classify(query)

        if intent == Intent.REFUSAL:
            return {
                "intent": intent.value,
                "response": "Diese Anfrage kann nicht "
                           "bearbeitet werden.",
                "cost": 0.0,
            }

        route = self.routes.get(intent, self.routes[Intent.SIMPLE])

        cached_response = semantic_cache.lookup(query)
        if cached_response[0]:
            return {
                "intent": intent.value,
                "response": cached_response[0],
                "cost": 0.0,
                "cached": True,
            }

        response = call_llm(
            route["model"],
            query,
            max_tokens=route["max_tokens"],
        )

        return {
            "intent": intent.value,
            "model": route["model"],
            "response": response,
            "cost": route["cost_per_call"],
            "cached": False,
        }

router = SemanticRouter()
result = router.route("Was ist der Status meiner Bestellung #12345?")
print(f"Intent: {result['intent']}, "
      f"Modell: {result.get('model', 'blocked')}, "
      f"Cost: EUR{result['cost']:.4f}")
\end{lstlisting}

\subsection{Router-Optimierung}

Der Router selbst kostet einen zusätzlichen LLM-Call (GPT-4o-mini). Bei hohem Volumen (10.000+ Requests/Tag) lohnt sich ein dedizierter Classifier:

\begin{itemize}[leftmargin=*]
    \item \textbf{Embedding-basiert}: Intent als Embedding + kNN-Klassifikation (10ms Latenz, kein LLM-Call)
    \item \textbf{kleiner Classifier}: Fine-getuntes MiniLM oder BERT für Intent-Klassifikation
    \item \textbf{Regel-basiert}: Regex-Patterns für einfache Intents (z.\,B. Bestellstatus, Öffnungszeiten)
\end{itemize}

\praxishinweis{Der Router selbst kostet einen LLM-Call. Ab 10.000 Requests/Tag lohnt sich ein dedizierter Intent-Classifier (MiniLM/BERT). Und immer Fallback definieren -- falsch klassifizierte Requests dürfen nicht einfach verschwinden.}

% ===========================================================%
% GUARDRAILS (VORWÄRTSVERWEIS AUF KAP. 16)
% ===========================================================%
\section{Guardrails -- Vorwärtsverweis auf Kap. 16 (Security)}

Guardrails sind Validierungsschichten vor und nach dem LLM-Call:

\begin{verbatim}
   [User Input] --> [Input Guard] --> [LLM] --> [Output Guard] --> [User]
                         |                           |
                    [PII Redact]              [Content Filter]
                    [Injection Check]           [Schema Validate]
                    [Length Limit]              [Hallucination Check]
\end{verbatim}

Die vollständige Guardrails-Architektur (Input Guard mit PII-Redaktion + Injection-Detection, Output Guard mit Content-Filter + Schema-Validation + Halluzinations-Check, Budget-Guard, Cost-Controller) findest du in \textbf{Kapitel 16 (Security)}.

Hier nur das Minimum für Produktion:
\begin{enumerate}
    \item \textbf{Input Guard}: PII-Redaktion + Injection-Check (Regex + LLM-Judge) + Length Limit
    \item \textbf{Output Guard}: Content-Filter (Safety) + Schema-Validation (Pydantic/Zod) + Halluzinations-Check (Faithfulness)
    \item \textbf{Budget Guard}: Cost/Token/Step Limits mit Circuit-Breaker
\end{enumerate}

\begin{lstlisting}[language=Python, caption=Output Guard mit Content-Filter und Schema-Validierung]
from typing import Optional


class OutputGuard:
    def __init__(self):
        self.blocked_categories = [
            "gewaltverherrlichend", "beleidigend",
            "diskriminierend", "illegal",
        ]
        self.max_length = 4096  # Zeichen

    def check_content_safety(self, text: str) -> Optional[str]:
        """LLM-basierter Content-Safety-Check."""
        prompt = (
            f"Pruefe den folgenden Text auf verbotene Inhalte. "
            f"Kategorien: {', '.join(self.blocked_categories)}.\n"
            f"Antworte NUR mit 'OK' oder 'BLOCKED: <Kategorie>'.\n\n"
            f"{text[:2000]}"
        )
        response = call_llm("gpt-4o-mini", prompt, temperature=0.0)
        if response.startswith("BLOCKED"):
            return response
        return None

    def validate_schema(self, text: str,
                        schema: dict) -> tuple[bool, str]:
        """Prueft, ob der Output einem JSON-Schema entspricht."""
        try:
            import json
            data = json.loads(text)
            for key, expected_type in schema.items():
                if key not in data:
                    return False, f"Missing key: {key}"
                if not isinstance(data[key], expected_type):
                    return False, (f"Wrong type for {key}: "
                                   f"{type(data[key])} != {expected_type}")
            return True, "OK"
        except json.JSONDecodeError as e:
            return False, f"Invalid JSON: {e}"

    def validate(self, llm_output: str,
                 schema: dict = None) -> dict:
        result = {
            "original_length": len(llm_output),
            "truncated": False,
            "blocked": False,
            "block_reason": None,
        }

        # Laengenbegrenzung
        if len(llm_output) > self.max_length:
            llm_output = llm_output[:self.max_length]
            result["truncated"] = True

        # Inhaltspruefung
        safety_issue = self.check_content_safety(llm_output)
        if safety_issue:
            result["blocked"] = True
            result["block_reason"] = safety_issue
            result["output"] = "[Output blocked by safety guard]"
            return result

        # Schema-Validierung
        if schema:
            valid, msg = self.validate_schema(llm_output, schema)
            result["schema_valid"] = valid
            result["schema_message"] = msg

        result["output"] = llm_output
        return result
\end{lstlisting}

\praxishinweis{Output Guard muss serverseitig laufen -- Client-Prüfung ist umgehbar. Safety-Check mit kleinem Modell (GPT-4o-mini) reicht für 90\% der Fälle. Nur für sensible Domänen ein zweites, strengeres Modell nachschalten.}

% ===========================================================%
% BUDGET CONTROL
% ===========================================================%
\section{Budget- und Rate-Limiting}

Ohne Kostenkontrolle kann ein LLM-System schnell explodieren. Ein fehlerhafter Agent, der in einer Schleife API-Calls auslöst, kostet innerhalb von Minuten hunderte Euro.

\begin{lstlisting}[language=Python, caption=Budget-Controller mit Rate-Limiting]
import time
from collections import defaultdict

class BudgetController:
    def __init__(self):
        self.daily_budget = 100.0    # EUR pro Tag (global)
        self.user_budget = 5.0       # EUR pro Tag (per User)
        self.rate_limit = 60         # Requests pro Minute (per User)
        self.user_costs = defaultdict(float)
        self.user_requests = defaultdict(list)
        self.total_cost = 0.0

    def check_request(self, user_id: str,
                      estimated_cost: float) -> dict:
        now = time.time()

        # 1. Globales Budget
        if self.total_cost + estimated_cost > self.daily_budget:
            return {"allowed": False,
                    "reason": "Globales Tagesbudget erschöpft"}

        # 2. User-Budget
        if (self.user_costs[user_id] + estimated_cost
                > self.user_budget):
            return {"allowed": False,
                    "reason": "User-Tagesbudget erschöpft"}

        # 3. Rate-Limit
        minute_ago = now - 60
        self.user_requests[user_id] = [
            t for t in self.user_requests[user_id]
            if t > minute_ago
        ]
        if len(self.user_requests[user_id]) >= self.rate_limit:
            return {"allowed": False,
                    "reason": "Rate-Limit erreicht (60/min)"}

        # Genehmigt
        self.user_requests[user_id].append(now)
        return {"allowed": True}

    def record_cost(self, user_id: str, cost: float):
        self.user_costs[user_id] += cost
        self.total_cost += cost

    def budget_status(self) -> dict:
        return {
            "total_spent": self.total_cost,
            "total_budget": self.daily_budget,
            "remaining": self.daily_budget - self.total_cost,
            "active_users": len(self.user_costs),
            "top_spenders": sorted(
                self.user_costs.items(),
                key=lambda x: -x[1],
            )[:5],
        }

# Nutzung im Request-Handler
budget = BudgetController()

def handle_request(user_id: str, query: str):
    estimated = estimate_cost(query)
    check = budget.check_request(user_id, estimated)

    if not check["allowed"]:
        return {"error": check["reason"]}

    response = call_llm("gpt-4o-mini", query)
    actual_cost = calculate_cost(response)
    budget.record_cost(user_id, actual_cost)

    return {"response": response, "cost": actual_cost}
\end{lstlisting}

\praxishinweis{Budget-Controller ohne Alarm nützen nichts. Bei 80\% Tagesbudget Slack-Benachrichtigung, bei 100\% Request-Queue pausieren. Ein fehlerhafter Agent in einer Schleife kostet in Minuten hunderte Euro.}

% ===========================================================%
% BEST PRACTICES
% ===========================================================%
\section{Best Practices}

\begin{itemize}[leftmargin=*]
    \item \textbf{Cache-Hierarchie}: Zuerst Response Cache prüfen, dann Semantic Cache, zuletzt LLM-Call. Reihenfolge = Kostenersparnis.
    \item \textbf{TTL anpassen}: FAQ-Responses: 24h Cache-Lebensdauer. Statusabfragen: 5 Minuten. Nicht-cachen: personalisierte oder sicherheitsrelevante Antworten.
    \item \textbf{Cache-Invalidierung}: Bei Prompt-Änderungen den gesamten Cache leeren (oder Version im Cache-Key führen).
    \item \textbf{Embedding-Modell}: Für Semantic Caching text-embedding-3-small (OpenAI) oder gte-small (Open Source). Grosse Modelle lohnen sich nicht, die Qualität steigt kaum.
    \item \textbf{Guardrails nicht im LLM}: Content-Filter mit Regex und kleinen Modellen (MiniLM) sind schneller und billiger als LLM-basierte Filter. Nur für komplexe Safety-Prüfungen GPT-4o-mini nutzen.
    \item \textbf{PII vor dem Cache}: Niemals PII im Cache speichern. Immer vor dem Caching redigieren.
    \item \textbf{Budget-Alarme}: Bei 80\,\% Tagesbudget: Slack/Teams-Benachrichtigung. Bei 100\,\%: Request-Queue pausieren.
\end{itemize}

% ===========================================================%
% ANTI-PATTERNS
% ===========================================================%
\section{Anti-Patterns}

\begin{itemize}[leftmargin=*]
    \item \textbf{Semantic Cache ohne Invalidierung}: Ein gecachtes Modell-Update vom Januar wird im Juni immer noch ausgeliefert. Immer TTL setzen.
    \item \textbf{Zu niedriger Threshold}: Ein Semantic Cache mit Threshold 0,8 liefert falsche Antworten -- die semantische Distanz ist zu gross. 0,92 ist der sicherere Start.
    \item \textbf{Keine PII-Prüfung vor Cache}: Persönliche Daten landen im Cache und werden an andere User ausgeliefert. Immer vor dem Caching redigieren.
    \item \textbf{Router ohne Fallback}: Wenn der Router eine Anfrage falsch klassifiziert (z.\,B. Code als Simple), gibt es keine Rückfallebene. Router-Confidence-Score mit Threshold implementieren.
    \item \textbf{Output Guard nur für den Browser}: Der Output Guard muss \textit{serverseitig} laufen -- nicht im Frontend. Client-seitige Prüfung ist umgehbar.
    \item \textbf{Budget ohne Monitoring}: Ein Budget-Controller nützt nichts, wenn ihn niemand überwacht. Immer Dashboards und Alarme für Budget-Metriken.
\end{itemize}

% ===========================================================%
% ZUSAMMENFASSUNG
% ===========================================================%
\begin{zusammenfassungEnv}
\begin{itemize}[leftmargin=*]
    \item Prompt Caching (KV-Cache) spart 50--90\,\% Prefill-Zeit und bis zu 50\,\% Input-Kosten
    \item Response Caching eliminiert wiederholte LLM-Calls bei exakt gleichen Inputs
    \item Semantic Caching (Embedding-Vergleich, Threshold $\sim$0,92) erfasst semantisch ähnliche Anfragen
    \item Der Semantic Router steuert Modell- und Pipeline-Wahl basierend auf Intent-Klassifikation
    \item Input Guard (PII, Injection) und Output Guard (Safety, Schema) sind Pflicht für Produktion
    \item Budget-Controller mit Rate-Limiting verhindert Kostenexplosion durch fehlerhafte Loops
\end{itemize}
\end{zusammenfassungEnv}




\par\mbox{}\par
\begin{merkeEnv}
Caching ist der effektivste Hebel zur Kostensenkung in LLM-Systemen -- 20--40\,\% aller API-Calls sind vermeidbar. Kombiniere Response Cache (exakt) mit Semantic Cache (ähnlich) und Prompt Caching (KV-Cache) für maximale Abdeckung. Guardrails sind keine Option, sondern eine Produktionsvoraussetzung.
\end{merkeEnv}

% ===========================================================%
% PRAXISPROJEKT
% ===========================================================%
\section{Praxisprojekt -- Caching-Strategie für einen Support-Chatbot}

\textbf{Ziel:} Baue eine Caching- und Guardrail-Strategie für einen Support-Chatbot mit 10.000 Requests/Tag.

\textbf{Aufgaben:}
\begin{enumerate}[leftmargin=*]
    \item Implementiere einen Semantic Cache mit Threshold 0,92 für FAQ-ähnliche Fragen
    \item Baue einen Input Guard, der E-Mails, Telefonnummern und IBANs redigiert
    \item Erstelle einen Semantic Router mit 4 Intents (Order, FAQ, Complaint, Other)
    \item Implementiere einen Budget-Controller mit 50~EUR Tageslimit und 20 Requests/min pro User
    \item Miss die Cache-Hit-Rate über 1.000 simulierte Requests
    \item Berechne die Kosteneinsparung durch den Cache
\end{enumerate}

\textbf{Erfolgskriterium:} Mindestens 25\,\% Cache-Hit-Rate, null PII-Leaks, null Budget-Überschreitungen.

% ===========================================================%
% RESSOURCEN
% ===========================================================%
\section{Ressourcen}

\begin{itemize}[leftmargin=*]
    \item \textbf{OpenAI Prompt Caching}: \href{https://platform.openai.com/docs/guides/prompt-caching}{platform.openai.com} -- Prefix-based KV-Cache
    \item \textbf{Anthropic Prompt Caching}: \href{https://docs.anthropic.com/en/docs/build-with-claude/prompt-caching}{docs.anthropic.com} -- cache\_control Header
    \item \textbf{vLLM Prefix Caching}: \href{https://docs.vllm.ai/en/latest/features/automatic_prefix_caching.html}{docs.vllm.ai} -- Automatic Prefix Caching
    \item \textbf{Semantic Router}: \href{https://github.com/aurelio-labs/semantic-router}{github.com/aurelio-labs/semantic-router} -- Open-Source Semantic Router
    \item \textbf{Guardrails AI}: \href{https://github.com/guardrails-ai/guardrails}{github.com/guardrails-ai/guardrails} -- Open-Source Guardrails Framework
    \item \textbf{NVIDIA NeMo Guardrails}: \href{https://github.com/NVIDIA/NeMo-Guardrails}{github.com/NVIDIA/NeMo-Guardrails} -- Enterprise Guardrails
\end{itemize}

\textbf{Nächster Schritt:} Mit Caching und Guardrails etabliert geht es weiter zur Inference-Optimierung für Self-Hosting -- vLLM, PagedAttention, Quantisierung, Speculative Decoding. Die Hardware-Ebene der Kostensenkung.
\endinput

\clearpage
